\section{置換積分の公式}\label{節:合成形の置換}

前章の一次式の積分では,内側の $ax+b$ の微分が定数 $a$ だったので,係数を直すだけで済んだ.では,内側が一般の関数 $g(x)$ のときはどうなるだろうか.ここで合成関数の微分を思い出そう.\,$F'(t)=f(t)$ であれば,
\[
 \frac{\dd}{\dd x}F(g(x))=f(g(x))g'(x)
\]
となる.この右辺の形が積分される式に現れたら,原始関数は左辺の $F(g(x))$ である.つまり,\,$f(g(x))g'(x)$ の積分は,\,$f$ の原始関数 $F$ に帰着する.これが置換積分の根拠である.

ここで,微分可能で,その導関数も連続である関数を\textbf{連続微分可能}という.\label{用語:連続微分可能}

\begin{theorem}{合成形の置換積分}
区間 $I$ で $g$ は連続微分可能とし,\,$g$ の値を含む区間で $f$ は連続とする.\,$F'=f$ ならば
\[
 \int f(g(x))g'(x)\dd x=F(g(x))+C.
\]
さらに $a,b\in I$ に対して
\[
 \int_a^b f(g(x))g'(x)\dd x
 =\int_{g(a)}^{g(b)}f(t)\dd t.
\]
この2つの式に,\,$g$ が単調であるという条件は要らない.
\end{theorem}

不定積分の式は,上で確かめた合成関数の微分そのものである.定積分の式も,基本定理によって両辺がともに
\[
 F(g(b))-F(g(a))
\]
となることから分かる.\,$g$ が途中で増加から減少へ変わっても,この証明は変わらない.後で逆関数を使う置換と比べるため,この点を覚えておこう.

実際の計算では,\,$t=g(x)$ と置き,\,$\dd t=g'(x)\dd x$ と書く.この記号は,\,$g'(x)$ と $\dd x$ を含む積を,\,$t$ による積分へ書き換えることを表している.微小な数を自由に割り算するという説明に頼る必要はなく,根拠は上の合成関数の微分にある.

\begin{bookpoint}
\textbf{POINT}\par
まとまり $g(x)$ を $t$ と置き,その微分 $g'(x)\dd x$ が残りの式にあることを確かめて $\dd t$ に換え,\,$t$ で積分して $x$ に戻す.
\end{bookpoint}

\begin{example}{合成関数の積分}
$\int x(x^2+1)^{20}\dd x$ を求めよ.
\end{example}

\paragraphbf{方針}
$t=x^2+1$ と置くと $\dd t=2x\dd x$ であり,残りの $x\dd x$ が $\frac12\dd t$ になる.

\paragraphbf{解答}
$\frac{(x^2+1)^{21}}{42}+C$

\paragraphbf{解法}
$t=x^2+1$ と置くと,\,$\dd t=2x\dd x$ だから,\,$x\dd x=\frac12\dd t$ とまとめて書き換えられる.したがって
\begin{align*}
 \int x(x^2+1)^{20}\dd x
 &=\frac12\int t^{20}\dd t\\
 &=\frac{t^{21}}{42}+C\\
 &=\frac{(x^2+1)^{21}}{42}+C
\end{align*}
となる.新しい変数 $t$ で積分した後,不定積分の答えは元の変数 $x$ に戻す.整式なので,実数全体で成り立つ.右辺を微分すると
\[
 \frac{21}{42}(x^2+1)^{20}\cdot2x
 =x(x^2+1)^{20}
\]
となり,元の式に戻る.

\paragraphbf{$x$ で割らずに置き換える}
解法では,\,$\dd x=\frac{\dd t}{2x}$ として $x$ で割ることはしなかった.\,$x\dd x$ 全体を $\frac12\dd t$ に換えたので,\,$x=0$ を除く必要がなく,答えは $x=0$ でも成り立つ.また $x^2+1$ は実数全体では単調でないが,原始関数の微分による確認には問題がない.定理に単調性の条件がないことと一致している.

\paragraphbf{まとまりの選び方}
新しい変数を選ぶときは,「見た目が複雑な部分」だけでなく,「その部分を微分したものが残りにあるか」を調べる.たとえば
\[
 \int\frac{\log x}{x}\dd x
\]
では,\,$t=\log x$ と置けば $\dd t=\frac1x\dd x$ であり,\,$\int t\dd t$ に変わる.したがって,\,$x>0$ で答えは $\frac12(\log x)^2+C$ となる.複雑に見える対数を変数にすると,同時に分母の $x$ もなくなるのだ.
